\subsection{Arquitetura Geral do Sistema}

Neste projeto, o cronômetro é modelado como uma máquina de estados finita cíclica $(Q, \Sigma, q_0, \delta, F)$, cujo comportamento evolui de forma determinística ao longo do tempo.
Tal que:

\begin{itemize}
    \item $Q = (H, M, S)$;
    \item $\Sigma = \{0,1\}$: 0 para incremento, 1 para decremento;
    \item $\delta : Q \times \Sigma \rightarrow Q$

onde:

\begin{align*}
H &\in [0, 23] \quad \text{representa as horas} \\
M &\in [0, 59] \quad \text{representa os minutos} \\
S &\in [0, 59] \quad \text{representa os segundos} \\
\end{align*}

O número total de estados possíveis do sistema é, portanto, finito e dado por:

\[
|Q| = 24 \cdot 60 \cdot 60 \cdot 2 = 172800
\]

\subsubsection{Estrutura interna do estado}

Existem duas principais formas de representar o estado interno do sistema, \textbf{representação binária direta} e \textbf{representação bcd}. 
Ambas são representações do sistema, porém com diferentes compromissos entre eficiência de informação e simplicidade de implementação, respectivamente.
Como demonstrado acima, existem $172800$ estados possíveis, o que requer um número mínimo $n$ de bits para representação, dado por:

\begin{align*}
n &= \left\lceil \log_2(24 \times 60 \times 60 \times 2) \right\rceil \\[4pt]
&= \left\lceil \log_2(24) + \log_2(60) + \log_2(60) + \log_2(2) \right\rceil \\[4pt]
&= \left\lceil (3 + \log_2 3) + 2(2 + \log_2 3 + \log_2 5) + 1 \right\rceil \\[4pt]
&= \left\lceil 3 + \log_2 3 + 4 + 2\log_2 3 + 2\log_2 5 + 1 \right\rceil \\[4pt]
&\approx \left\lceil 8 + 3(1.585) + 2(2.322) \right\rceil \\[4pt]
&= \left\lceil 17.399 \right\rceil \\[4pt]
&= 18
\end{align*}

Por outro lado, a representação BCD utiliza 4 bits para cada dígito decimal, totalizando 6 dígitos (2 para horas, 2 para minutos e 2 para segundos), resultando em $6 \cdot 4 = 24$ bits.
Entretanto, note que, sob as restrições deste sistema, a representação BCD apresenta algumas vantagens em detrimento da representação binária direta.

Primeiramente, além de ser mais intuitiva para a representação de valores decimais, a estrutura modular do BCD facilita a implementação de contadores independentes para cada dígito, simplificando a lógica de incremento e decremento. 

Imagine, por exemplo, que fosse utilizada a representação binária direta, onde o estado $S$ é representado por um registrador de 18 bits. Para projetar uma função de transição LUT (), global, que realize o incremento ou decremento, seria necessário lidar com a complexidade de manipular um número binário de 18 bits, o que envolve
tabelas verdade com $2^{18} = 262144$ entradas, tornando a implementação impraticável. 
Por mais ue sistemas reais utilizem tecnicas combinacionais para reduzir a complexidade, a implementação em representação binária direta ainda ainda implicaria em uma função de transição de alta dimensionalidade, dificultando sua decomposição modular.

Em contraste, com a representação BCD, cada dígito pode ser tratado como um contador independente, cada um com um número pequeno de bits.
Além disso, dado os intervalos específicos de cada unidade de contagem (0-23 para horas, 0-59 para minutos e segundos), a representação BCD pode ser otimizada para reduzir a quantidade de bits necessários para armazenar o estado, eliminando a necessidade de bits adicionais para representar estados fora desses intervalos. Nesse contexto:

\[
H = (H_{d}, H_{u}), \quad M = (M_{d}, M_{u}), \quad S = (S_{d}, S_{u})
\]

com:

\[
H_{d} \in [0,2], \quad H_{u} \in [0,9], \quad M_{d}, S_{d} \in [0,5], \quad M_{u}, S_{u} \in [0,9]
\]

Assim, a representação BCD para o estado $Q$ pode ser organizada em um vetor de bits $D$ de tamanho $N = 21$, onde cada bit é alocado para representar os dígitos BCD correspondentes, conforme a seguinte estrutura: 

\begin{center}
\begin{tabular}{c|c}
Bit & Função \\
\hline
0 & Controle de direção (0 para incremento, 1 para decremento)\\
$(0,2]$ & $H_d$ \\
$(2,6]$ & $H_u$ \\
$(6,9]$ & $M_d$ \\
$(9,13]$ & $M_u$  \\
$(13,16]$ & $S_d$ \\
$(16,20]$ & $S_u$ \\
\end{tabular}
\end{center}

Dessa forma, a função de transição $\delta$ pode ser estruturada como um conjunto de subfunções independentes associadas a cada componente do estado, isto é, horas, minutos e segundos. Essa decomposição explora a estrutura hierárquica do sistema, reduzindo o acoplamento entre variáveis globais e permitindo que cada submódulo opere sobre um subconjunto restrito de bits.

Como consequência, a complexidade da implementação não é dominada pela cardinalidade do espaço global de estados, mas sim pela complexidade local de cada unidade de atualização e pelas condições de propagação entre dígitos (carry/borrow). Isso viabiliza a construção de uma lógica combinacional modular, na qual cada bloco de transição pode ser projetado e otimizado de forma independente.

\subsubsection{Função de transição}

Como abordado na fundamentação acerca \nameref{sec:Máquina de Estados Finita}, a dinâmica do sistema é descrita por uma função de transição determinística $\delta$, que mapeia o estado atual do sistema para seu estado sucessor a cada evento de clock.

Seja $Q_t = (H_t, M_t, S_t)$ o estado do sistema no instante $t$. A evolução temporal é definida por:

\[
Q_{t+1} = \delta(Q_t)
\]

No contexto deste projeto, a função de transição é induzida por uma operação de incremento ou decremento, definida pelo parâmetro $\sigma_t \in \Sigma$, unitário de tempo, podendo ser escrita como:

\[
\delta(H_t, M_t, S_t; \sigma_t) = (H_{t+1}, M_{t+1}, S_{t+1})
\]

A atualização do estado não ocorre de forma uniforme sobre todos os componentes, mas segue uma estrutura hierárquica com propagação condicional entre dígitos.

Em particular, define-se a transição dos segundos como:

\[
(S_{t+1}, c_1) = \delta_S(S_t; \sigma_t)
\]

onde $c_1 \in \{0,1\}$ representa o carry de propagação para a unidade de minutos.

Analogamente, a transição dos minutos é dada por:

\[
(M_{t+1}, c_2) = \delta_M(M_t, c_1; \sigma_t)
\]

e a transição das horas por:

\[
(H_{t+1}) = \delta_H(H_t, c_2; \sigma_t)
\]

Dessa forma, a função global $\delta$ pode ser expressa como a composição hierárquica:

\[
\delta = \delta_H \circ \delta_M \circ \delta_S
\]

